% ID: SEL-D0019-OB-E01
% ESTADO: cotejada. Segunda pasada visual completa y borrador compilado revisados.
% Fuente: Selectividad/matematicas_aplicadas/reserva a/Examen-Matemáticas Aplicadas a las Ciencias Sociales.pdf
% SHA-256: 6fed0384ff6fcbaaa7098ab3f35ea5f881d1280502f050ada3a7f13af528392e
% Página: 2; opción B; ejercicio original 1.

Se consideran las matrices
$A=\begin{pmatrix}-1&2\\-3&4\end{pmatrix}$,
$B=\begin{pmatrix}-1&2&1\\3&0&2\end{pmatrix}$ y
$C=\begin{pmatrix}3&0&1\\2&-1&-1\end{pmatrix}$.
\begin{enumerate}[label=\alph*)]
\item \textbf{(0.5 puntos)} Razone qué dimensiones deben tener las matrices $P$ y $Q$ para que los productos $(A\cdot P\cdot B^t)$ y $(Q\cdot A\cdot C)$ den como resultado una matriz cuadrada.
\item \textbf{(2 puntos)} Resuelva la ecuación matricial $A\cdot X-2B\cdot C^t=A^2$
\end{enumerate}
